\section{Evaluation}
\label{sec:evaluation}

\subsection{Evaluation Questions}

The evaluation should answer four questions:
\begin{itemize}
\item \textbf{RQ1:} Does noise-aware training substantially improve performance over direct noisy training?
\item \textbf{RQ2:} Do the gains hold on safety-relevant metrics such as emergency recall and emergency F1?
\item \textbf{RQ3:} Which training ideas are already well supported, and which are still exploratory?
\item \textbf{RQ4:} Does the prototype remain viable under onboard latency and real-world acoustic conditions?
\end{itemize}

\subsection{Current Completed Benchmark Evidence}

Table~\ref{tab:current_noisy_results} consolidates the strongest completed noisy-set results currently visible in the repository. All runs below use evaluation SNR $=-10$ dB on a balanced three-class test set. The table is intentionally simple because the main writing need right now is to anchor the narrative with trustworthy evidence.

\begin{table}[t]
    \centering
    \caption{Current completed noisy-set results from the repository.}
    \label{tab:current_noisy_results}
    \begin{tabular}{lccc}
        \toprule
        Variant & Acc. & Emerg. R & Emerg. F1 \\
        \midrule
        Direct noisy baseline (`E0`) & 0.72 & 0.47 & 0.57 \\
        Stats+KD draft (`E2\_newTS`) & 0.79 & 0.66 & 0.73 \\
        CE+logits KD (`branch\_trial`) & 0.84 & 0.69 & 0.79 \\
        Preprocess+embed KD (`preprocess\_ext`) & 0.88 & 0.79 & 0.87 \\
        \bottomrule
    \end{tabular}
\end{table}

Two conclusions are already defensible from this table. First, transfer-based and preprocessing-aware variants consistently outperform the direct noisy baseline. Second, the paper should prioritize the `preprocess\_ext` and `branch\_trial` families because they currently offer the clearest evidence. By contrast, the newer stats-branch line is promising but not yet the strongest completed result.

\subsection{What the Current Results Say}

The current evidence supports a narrow but important claim: robust emergency intent recognition under rotor-like noise depends on \emph{how} supervision is transferred and regularized, not merely on adding more noisy data. The direct baseline misses over half of emergency commands on the noisy test set, whereas the best completed configuration raises emergency recall by $32$ percentage points.

The same results also constrain the paper story. Some auxiliary ideas can improve performance relative to the baseline, but not all of them beat the strongest preprocessing and distillation configuration. This is why the main paper should be written around the best-supported training recipe instead of promising every exploratory branch as a co-equal contribution.

\subsection{Still-Missing Evaluation for SenSys}

The benchmark evidence above is necessary, but it is not enough for submission. The next writing round must add:
\begin{itemize}
\item onboard latency and memory measurements;
\item real-world rotor-noise tests with a fixed protocol;
\item false-trigger analysis for the \emph{unknown} and emergency paths;
\item at least one ablation that isolates preprocessing, augmentation, and distillation contributions.
\end{itemize}

These items are not optional polish. They are the difference between a strong speech-model draft and a systems paper with a credible deployment claim.
